\documentclass[10pt]{article}

\usepackage[utf8]{inputenc}

\usepackage[T1]{fontenc}

\usepackage{amsmath}

\usepackage{amsfonts}

\usepackage{amssymb}

\usepackage[version=4]{mhchem}

\usepackage{stmaryrd}

\usepackage{graphicx}

\usepackage[export]{adjustbox}

\graphicspath{ {./images/} }



\begin{document}

любые гомоморфизмы $\varphi_{i}: G_{i} \rightarrow H$ групп $G_{i}$ в любую групіу $H$ продолжаются до гомоморфизма $\varphi: G \rightarrow H$. Для обозначения С. п. используется знак $*$, напр.:



\[

G=\prod_{i \in I}^{*} G_{i} \quad \text { и } \quad G=G_{1} * G_{2} * \ldots * G_{k}

\]



в случае конечного множества I. Каждый не равный единице элемент С. П. $G$ единственным образом выражается в виде несократимого слова $v=g_{i_{1}} g_{i_{2}} . g_{i_{n}}$, где $g_{i_{j}} \in G_{i_{j}}, g_{i_{j}} \neq 1$, и при любом $j=1,2, \ldots, n-1$, $i_{j} \neq i_{j+1}$. Конструкция С. п. является важной в изучении групп, заданных множеством порождающих әлементов и определяющих соотношений. В этих терминах оно может быть определено следующим образом. Пусть каждая группа $G_{i}$ задана множествами $X_{i}$ порождающих и $\Phi_{i}$ определяющих соотношений, причем $X_{i} \cap X_{j}=\phi$, если $i \neq j$. Тогда группа $G$, заданная множеством $X_{i}=\bigcup_{i \in I} X_{i}$ порождающих и множеством $\Phi=$ $=\bigcup_{i \in I} \Phi_{i}$ определяющих соотношений, будет С. п.



\begin{center}

\includegraphics[max width=\textwidth]{2023_07_08_ebbc368ece8aee55b542g-1}

\end{center}



Всякая подгруппа С. П. $G$ сама разлагается в С. п. своих подгрупп, из к-рых нек-рые являются бесконечными циклическими, а каждая из других сопряжена с нек-рой подгруппой какой-либо группы $G_{i}$, входящей в свободное разложение группы $G$ (т е о р е м а К уp o III a).



[2] Лum.: [1] К у $\mathrm{p}$ о т А. Г., Теория групп, 3 изд., М., 1967; [2] м а г н у с В., К а p p a с А., С о ли т ә p Д.,, Комбинаторная теория групп, пер. с англ., М., 1974 .



СВОБОДНЫЙ ВЕКТОР - см. Вектор. А. Л. Шмелькин.



СВОБОДНЫЙ ГРУППОИД - свободная алгебра многообразия всех группоидов. С. г. с множеством $X$ свободных образующих совпадает с группоидом слов, если под словом понимать любую упорядоченную систему элементов из $X$ с любыми повторениями, причем в этой системе задано распределение скобок (каждый символ считается взятым в скобки, а затем скобки расставлены так, что каждый раз перемножаются лишь две скобки). Произведением слов $(a)$ и (b) считается слово $((a)(b))$. О. А. Иванова.



СВОБОДНЫЙ МОДУЛЬ - свободный объект (свободная алгебра) в многообразии модулей над фиксированным кольцом $R$. Если $R$ - ассоциативное кольцо с единицей, то С. м. - это модуль, обладающий базисом, т. е. линейно независимой системой порождающих. Мощность базиса С. м. наз. ег'о р а н г о м. Ранг не всегда определен однозначно, т. е. существуют кольца, над к-рыми С. м. может обладать двумя базисами, состоящими из различного числа элементов. Это равносильно существованию над кольцом $R$ двух прямоугольных матриц $A$ и $B$, для к-рых



\[

A B=I_{m}, B A=I_{n}, m \neq n,

\]



где $I_{m}$ и $I_{n}$ - единичные матрицы порядков $m$ и $n$ соответственно. Неоднозначность, однако, имеет место лишь для конечных базисов, если же ранг С. м. бесконечен, то все базисы С. м. равномощны. Кроме того, над кольцами, допускающими гомоморфизм в тело (в частности, над коммутативными кольцами), ранг С. м. определен всегда однозначно.



Кольцо $R$, рассматриваемое как левый модуль над самим собой, является С. м. ранга один. Всякий левый С. м. является прямой суммой С. м. ранга один. Јюбой модуль $M$ представим как фактормодуль $F_{0} / H_{0}$ нек-рого С. м. $F_{0}$. Подмодуль $H_{0}$ в свою очередь представим как фактормодуль $F_{1} / H_{1}$ С. м. $F_{1}$. Продолжая этот процесс, получают точную последовательность:



\[

\cdots \rightarrow F_{2} \rightarrow F_{1} \rightarrow F_{0} \rightarrow M \rightarrow 0

\]



к-рая наз. с в о бод н о ӥ р е зо ль в ен т о й модуля $M$. Тела могут быть охарактеризованы как кольца, над к-рыми все модули свободны. Над областью главных идеалов подмодуль С. м. свободен. Близкими к С. М. являются проективные модули и плоские модули. Лит.: [1] К о н П., Свободные кольца и их связи, пер. с англ., М., 1975; [2] М а к л е й н С., Гомология, пер. с англ., СВОДИМОСТИ АКСИОМА - аксиома, добавленная Б. Расселом (B. Russell) к его разветвленной теории типов с целью избежать расслоения понятий (см. Henpeдикативное определение). В разветвленной теории типов множества данного типа разделяются на порядки. Так, вместо понятия множества натуральных чисел появляется понятие множества натуральных чисел данного порядка. При этом множество натуральных чисел, определяемое формулами без использования каких-либо множеств, принадлежит первому порядку. Если в определении используется совокупность множеств первого порядка, а совокупности множеств высшего порядка не используются, то определяемое множество принадлежит второму порядку, и т. д. Напр., если $S-$ семейство множеств, состоящее из множеств одного и того же порядка, то множество



\[

M=\{x \mid \exists y(y \in S \wedge x \in y)\}

\]



должно принадлежать следующему порядку, т. к. в его определении содержится квантор по множествам данного порядка. С. а. утверждает, что для каждого множества существует равнообъемное ему (т. е. состоящее из тех же самых әлементов) множество первого порядка. Таким образом, С. а. фактически сводит разветвленную теорию типов $к$ простой теории типов.



Лuт.: [1] $\Gamma$ и ль б е р т Д., Аккерма н В., Основы теоретической логики, пер. с нем., М., 1947.



СВЯЗАННАЯ ПЕРЕМЕННАЯ, с в я з а н н о е в хо жд ени е пе р е менн ой,- тип вхождения переменной в языковое выражение. Точное определение для каждого формализованного языка - свое и зависит от правил образования этого языка. Вместо С. п. нельзя подставлять объекты. Такая подстановка приводит к бессмысленным выражениям. Но замена С. П. всюду, где она встречается, на новую для данного выражения переменную приводит к выражению с тем же самым смыслом. Напр., в выражениях



\[

\int f(x, y) d x,\{x \mid f(x, y)=0\}

\]



переменная $x$ является связанной. Подстановка вместо $x$ какого-нибудь числа приводит к бессмысленным выражениям В то же время, написав всюду вместо $x$ например $z$, получают выражения, обозначающие те же самые сущности.



С. п. всегда возникают при применении к нек-рому выражению $e$ со свободными вхождениями переменной $x$ какого-нибудь оператора с операторной переменной $x$ (см. Свободная переменная). В получившемся выражении все вхождения переменной $x$ в $e$, бывпие свободными, становятся связанными. Ниже указаны нек-рые наиболее употребительные операторы (помимо уже использованных операторов $\int \ldots . . d x$ и $\left.\{x \mid \ldots\}\right)$, в к-рых $x$ является операторной переменной:



$\forall x(\ldots), \exists x(\ldots)$ - кванторы общности и существования;



$\int_{\Sigma}^{\cdots} \ldots d x$ - определенный интеграл по $x$;



$\Sigma_{x} \ldots-$ сумма по $x$



$\lambda x(\ldots)$ - функция от $x$, значение к-роӥ в точке $x$ равно ....



Вместо многоточий можно подставлять огределенные языковые выражения.



В реальных (не формализованных) мачематич. текстах возможно неоднозначное употребление одних и тех же выражений, в связи с чем выделение С. П. в данном





\end{document}